\documentclass[11pt]{article}

\usepackage[a4paper,margin=1in]{geometry}
\usepackage{amsmath,amssymb,amsthm,mathtools,bm}
\usepackage{booktabs,multirow,array}
\usepackage{enumitem}
\usepackage{hyperref}

\newtheorem{theorem}{Theorem}
\newtheorem{proposition}{Proposition}
\newtheorem{corollary}{Corollary}
\newtheorem{remark}{Remark}
\newtheorem{definition}{Definition}

\newcommand{\R}{\mathbb{R}}
\newcommand{\norm}[1]{\left\|#1\right\|}
\newcommand{\fro}[1]{\left\|#1\right\|_F}
\newcommand{\op}[1]{\left\|#1\right\|_2}
\newcommand{\tr}{\mathrm{tr}}
\newcommand{\diag}{\mathrm{diag}}
\newcommand{\Span}{\mathrm{span}}
\newcommand{\rank}{\mathrm{rank}}
\newcommand{\sinTheta}{\mathrm{sin}\,\Theta}

\title{Soft Mixing Regularization for Spectrally Structured PEFT:\\
Definitions, Perturbation Bounds, and Ablation Protocol}
\author{}
\date{}

\begin{document}
\maketitle

\section{Setup}

Let the pretrained weight matrix be
\[
W^{\mathrm{pre}} = U \Sigma V^\top
=
\begin{bmatrix} U_R & U_r \end{bmatrix}
\begin{bmatrix}
\Sigma_R & 0\\
0 & \Sigma_r
\end{bmatrix}
\begin{bmatrix} V_R & V_r \end{bmatrix}^\top,
\]
where:
\begin{itemize}[leftmargin=2em]
    \item \(U_R \in \R^{m\times R}\), \(V_R \in \R^{n\times R}\) are the leading singular vectors,
    \item \(U_r \in \R^{m\times r}\), \(V_r \in \R^{n\times r}\) are the complementary/tail singular vectors,
    \item \(U=[U_R\ U_r]\) and \(V=[V_R\ V_r]\) are orthogonal.
\end{itemize}

Consider the structured update
\[
\Delta W = E\,U_r\,\Sigma_r' \,V_r^\top D,
\]
where:
\begin{itemize}[leftmargin=2em]
    \item \(E\in\R^{m\times m}\) and \(D\in\R^{n\times n}\) are diagonal trainable scaling matrices,
    \item \(\Sigma_r' \in \R^{r\times r}\) is trainable (typically diagonal, but the results below only require it to be an \(r\times r\) matrix).
\end{itemize}

The adapted weight is
\[
\widetilde W = W^{\mathrm{pre}} + \Delta W.
\]

\section{Left and Right Mixing Terms}

The key quantities governing whether the update leaks into the preserved leading singular structure are the singular-basis cross terms induced by \(E\) and \(D\).

\begin{definition}[Left and right mixing matrices]
Define
\[
M_E := U_R^\top E U_r \in \R^{R\times r},
\qquad
M_D := V_R^\top D V_r \in \R^{R\times r}.
\]
Their operator-norm magnitudes are
\[
\mu_E := \op{M_E},
\qquad
\nu_D := \op{M_D}.
\]
\end{definition}

\begin{remark}
The condition \(\mu_E=0\) means that \(E\) does not map any tail left singular direction into the preserved leading left singular subspace. Likewise, \(\nu_D=0\) means that \(D\) does not mix the tail right singular directions into the preserved leading right singular subspace.
\end{remark}

In practice, it is often easier to optimize a Frobenius penalty:
\[
\mathcal{R}_{\mathrm{mix}}
=
\lambda_E^{\mathrm{mix}} \fro{M_E}^2
+
\lambda_D^{\mathrm{mix}} \fro{M_D}^2,
\]
since
\[
\mu_E = \op{M_E} \le \fro{M_E},
\qquad
\nu_D = \op{M_D} \le \fro{M_D}.
\]
Thus minimizing \(\mathcal{R}_{\mathrm{mix}}\) also controls \(\mu_E\) and \(\nu_D\).

\section{Blockwise Perturbation Analysis}

Write the perturbation in the singular coordinates of \(W^{\mathrm{pre}}\):
\[
C := U^\top \Delta W V.
\]
Partition \(C\) according to the \(R/r\) splitting:
\[
C=
\begin{bmatrix}
C_{11} & C_{12}\\
C_{21} & C_{22}
\end{bmatrix}.
\]
A direct calculation gives
\begin{equation}
\label{eq:block-form}
C
=
\begin{bmatrix}
U_R^\top E U_r\, \Sigma_r' \,V_r^\top D V_R
&
U_R^\top E U_r\, \Sigma_r' \,V_r^\top D V_r
\\[0.75em]
U_r^\top E U_r\, \Sigma_r' \,V_r^\top D V_R
&
U_r^\top E U_r\, \Sigma_r' \,V_r^\top D V_r
\end{bmatrix}.
\end{equation}

The meanings of these blocks are:
\begin{itemize}[leftmargin=2em]
    \item \(C_{11}\): direct perturbation of the preserved leading block,
    \item \(C_{12}\): coupling from the preserved left block to the tail right block,
    \item \(C_{21}\): coupling from the tail left block to the preserved right block,
    \item \(C_{22}\): perturbation internal to the tail block.
\end{itemize}

\begin{proposition}[Blockwise perturbation bounds]
\label{prop:block-bounds}
Let \(\Delta W = E U_r \Sigma_r' V_r^\top D\) and let \(C = U^\top \Delta W V\) be partitioned as in \eqref{eq:block-form}. Then:
\begin{align}
\op{C_{11}}
&\le
\mu_E\, \op{\Sigma_r'}\, \nu_D,
\label{eq:C11-bound}
\\
\op{C_{12}}
&\le
\mu_E\, \op{\Sigma_r'}\, \op{D},
\label{eq:C12-bound}
\\
\op{C_{21}}
&\le
\op{E}\, \op{\Sigma_r'}\, \nu_D,
\label{eq:C21-bound}
\\
\op{C_{22}}
&\le
\op{E}\, \op{\Sigma_r'}\, \op{D}.
\label{eq:C22-bound}
\end{align}
Consequently,
\begin{equation}
\label{eq:global-perturb-bound}
\op{\Delta W}
\le
\op{E}\,\op{\Sigma_r'}\,\op{D}.
\end{equation}
\end{proposition}

\begin{proof}
From \eqref{eq:block-form},
\[
C_{11}
=
(U_R^\top E U_r)\,\Sigma_r'\,(V_r^\top D V_R).
\]
Using submultiplicativity of the operator norm,
\[
\op{C_{11}}
\le
\op{U_R^\top E U_r}\,\op{\Sigma_r'}\,\op{V_r^\top D V_R}
=
\mu_E\,\op{\Sigma_r'}\,\nu_D.
\]
This proves \eqref{eq:C11-bound}.

Similarly,
\[
C_{12}
=
(U_R^\top E U_r)\,\Sigma_r'\,(V_r^\top D V_r),
\]
so
\[
\op{C_{12}}
\le
\op{U_R^\top E U_r}\,\op{\Sigma_r'}\,\op{V_r^\top D V_r}.
\]
Since \(V_r\) has orthonormal columns,
\[
\op{V_r^\top D V_r}\le \op{D},
\]
hence \eqref{eq:C12-bound} follows. The proof of \eqref{eq:C21-bound} is identical:
\[
\op{C_{21}}
\le
\op{U_r^\top E U_r}\,\op{\Sigma_r'}\,\op{V_r^\top D V_R}
\le
\op{E}\,\op{\Sigma_r'}\,\nu_D.
\]
Finally,
\[
\op{C_{22}}
\le
\op{U_r^\top E U_r}\,\op{\Sigma_r'}\,\op{V_r^\top D V_r}
\le
\op{E}\,\op{\Sigma_r'}\,\op{D},
\]
which proves \eqref{eq:C22-bound}.

For the global bound, since \(U\) and \(V\) are orthogonal,
\[
\op{\Delta W}=\op{C}.
\]
Also,
\[
\Delta W = E U_r \Sigma_r' V_r^\top D,
\]
so by submultiplicativity and orthonormality of \(U_r,V_r\),
\[
\op{\Delta W}
\le
\op{E}\,\op{U_r}\,\op{\Sigma_r'}\,\op{V_r^\top}\,\op{D}
=
\op{E}\,\op{\Sigma_r'}\,\op{D}.
\]
\end{proof}


\begin{proposition}[Tail decomposition, off-tail bound, and exact tail confinement]
\label{prop:tail-confinement}
Let
\[
\Delta W = E U_r \Sigma_r' V_r^\top D,
\qquad
C := U^\top \Delta W V
=
\begin{bmatrix}
C_{11} & C_{12}\\
C_{21} & C_{22}
\end{bmatrix},
\]
where the block partition is with respect to the pretrained singular splitting
\(
U=[U_R\ U_r]
\)
and
\(
V=[V_R\ V_r].
\)
Define the pure tail part and the off-tail part by
\[
C_{\mathrm{tail}}
:=
\begin{bmatrix}
0 & 0\\
0 & C_{22}
\end{bmatrix},
\qquad
C_{\mathrm{off}}
:=
C-C_{\mathrm{tail}}
=
\begin{bmatrix}
C_{11} & C_{12}\\
C_{21} & 0
\end{bmatrix}.
\]
Then:
\begin{enumerate}
    \item[(i)] \textbf{(Exact decomposition)} One has
    \[
    C = C_{\mathrm{tail}} + C_{\mathrm{off}},
    \]
    with
    \[
    C_{11} = (U_R^\top E U_r)\,\Sigma_r'\,(V_r^\top D V_R),
    \]
    \[
    C_{12} = (U_R^\top E U_r)\,\Sigma_r'\,(V_r^\top D V_r),
    \]
    \[
    C_{21} = (U_r^\top E U_r)\,\Sigma_r'\,(V_r^\top D V_R),
    \]
    \[
    C_{22} = (U_r^\top E U_r)\,\Sigma_r'\,(V_r^\top D V_r).
    \]

    \item[(ii)] \textbf{(Off-tail norm bound)} The off-tail part satisfies
    \[
    \|C-C_{\mathrm{tail}}\|_2
    =
    \|C_{\mathrm{off}}\|_2
    \le
    \left\|
    \begin{bmatrix}
    \|C_{11}\|_2 & \|C_{12}\|_2\\
    \|C_{21}\|_2 & 0
    \end{bmatrix}
    \right\|_2
    \le
    \sqrt{
    \|C_{11}\|_2^2+\|C_{12}\|_2^2+\|C_{21}\|_2^2
    }.
    \]

    \item[(iii)] \textbf{(Mixing-controlled off-tail bound)} If
    \[
    \mu_E := \|U_R^\top E U_r\|_2,
    \qquad
    \nu_D := \|V_R^\top D V_r\|_2,
    \]
    then
    \[
    \|C-C_{\mathrm{tail}}\|_2
    \le
    \|\Sigma_r'\|_2
    \sqrt{
    (\mu_E\nu_D)^2
    +
    (\mu_E\|D\|_2)^2
    +
    (\|E\|_2\nu_D)^2
    }.
    \]

    \item[(iv)] \textbf{(Exact tail confinement)} If
    \[
    \mu_E = 0
    \qquad\text{and}\qquad
    \nu_D = 0,
    \]
    then
    \[
    C_{11}=0,\qquad C_{12}=0,\qquad C_{21}=0,
    \]
    and therefore
    \[
    C=C_{\mathrm{tail}}
    =
    \begin{bmatrix}
    0 & 0\\
    0 & C_{22}
    \end{bmatrix}.
    \]
    Equivalently, in pretrained singular coordinates, the perturbation is confined entirely to the tail block.
\end{enumerate}
\end{proposition}

\begin{proof}
Part (i) follows by direct multiplication:
\[
C
=
U^\top (E U_r \Sigma_r' V_r^\top D) V.
\]
Partitioning with respect to \(U=[U_R\ U_r]\) and \(V=[V_R\ V_r]\) yields the stated expressions for \(C_{11},C_{12},C_{21},C_{22}\), hence also
\(
C = C_{\mathrm{tail}} + C_{\mathrm{off}}.
\)

For part (ii), define
\[
a:=\|C_{11}\|_2,\qquad
b:=\|C_{12}\|_2,\qquad
c:=\|C_{21}\|_2.
\]
Then
\[
C_{\mathrm{off}}
=
\begin{bmatrix}
C_{11} & C_{12}\\
C_{21} & 0
\end{bmatrix}.
\]
For any block vector \(x=(x_1,x_2)\),
\[
C_{\mathrm{off}}x
=
\begin{bmatrix}
C_{11}x_1 + C_{12}x_2\\
C_{21}x_1
\end{bmatrix},
\]
so
\[
\|C_{\mathrm{off}}x\|_2
\le
\left\|
\begin{bmatrix}
a & b\\
c & 0
\end{bmatrix}
\begin{bmatrix}
\|x_1\|_2\\
\|x_2\|_2
\end{bmatrix}
\right\|_2.
\]
Taking the supremum over \(\|x\|_2=1\) gives
\[
\|C_{\mathrm{off}}\|_2
\le
\left\|
\begin{bmatrix}
a & b\\
c & 0
\end{bmatrix}
\right\|_2.
\]
Finally, the operator norm is bounded by the Frobenius norm:
\[
\left\|
\begin{bmatrix}
a & b\\
c & 0
\end{bmatrix}
\right\|_2
\le
\left\|
\begin{bmatrix}
a & b\\
c & 0
\end{bmatrix}
\right\|_F
=
\sqrt{a^2+b^2+c^2},
\]
which proves part (ii).

For part (iii), apply Proposition~\ref{prop:block-bounds}:
\[
\|C_{11}\|_2 \le \mu_E \|\Sigma_r'\|_2 \nu_D,
\]
\[
\|C_{12}\|_2 \le \mu_E \|\Sigma_r'\|_2 \|D\|_2,
\]
\[
\|C_{21}\|_2 \le \|E\|_2 \|\Sigma_r'\|_2 \nu_D.
\]
Substituting these into part (ii) yields
\[
\|C-C_{\mathrm{tail}}\|_2
\le
\|\Sigma_r'\|_2
\sqrt{
(\mu_E\nu_D)^2
+
(\mu_E\|D\|_2)^2
+
(\|E\|_2\nu_D)^2
}.
\]

For part (iv), if \(\mu_E=0\), then
\[
U_R^\top E U_r = 0.
\]
If \(\nu_D=0\), then
\[
V_R^\top D V_r = 0.
\]
Because \(D\) is diagonal (hence symmetric over \(\mathbb{R}\)),
\[
V_r^\top D V_R = (V_R^\top D V_r)^\top = 0.
\]
Substituting into the formulas in part (i) yields
\[
C_{11}=0,\qquad C_{12}=0,\qquad C_{21}=0.
\]
Therefore
\[
C=C_{\mathrm{tail}}
=
\begin{bmatrix}
0 & 0\\
0 & C_{22}
\end{bmatrix}.
\]
\end{proof}

\begin{corollary}[Asymptotic tail confinement under soft mixing control]
\label{cor:asymptotic-tail}
Assume that along training (or along a sequence of iterates) the mixing quantities satisfy
\[
\mu_E = \|U_R^\top E U_r\|_2 \to 0,
\qquad
\nu_D = \|V_R^\top D V_r\|_2 \to 0,
\]
and that \(\|E\|_2\), \(\|D\|_2\), and \(\|\Sigma_r'\|_2\) remain bounded.
Then
\[
\|C-C_{\mathrm{tail}}\|_2 \to 0.
\]
Equivalently, the perturbation converges in operator norm to a pure tail-block perturbation:
\[
C \to
\begin{bmatrix}
0 & 0\\
0 & C_{22}
\end{bmatrix}.
\]
\end{corollary}

\begin{proof}
This follows immediately from Proposition~\ref{prop:tail-confinement}(iii).
\end{proof}

\begin{remark}
Proposition~\ref{prop:tail-confinement} and Corollary~\ref{cor:asymptotic-tail}
make precise that shrinking the mixing terms does \emph{not} force the full perturbation \(C\) to vanish.
Rather, it forces only the \emph{off-tail part}
\(
C-C_{\mathrm{tail}}
\)
to vanish, so that the perturbation becomes confined to the tail block.
\end{remark}


\section{Soft Mixing Regularization and Approximate Preservation}

We now turn the above blockwise analysis into a theorem that can be used in the paper.

% \begin{theorem}[Soft mixing regularization yields approximate block preservation]
% \label{thm:soft-mix}
% Consider the training objective
% \[
% \mathcal{L}_{\mathrm{total}}
% =
% \mathcal{L}_{\mathrm{task}}
% +
% \lambda_E^{\mathrm{mix}}\fro{U_R^\top E U_r}^2
% +
% \lambda_D^{\mathrm{mix}}\fro{V_R^\top D V_r}^2.
% \]
% Suppose that at some iterate (or at a local optimum) we have
% \[
% \fro{U_R^\top E U_r}\le \varepsilon_E,
% \qquad
% \fro{V_R^\top D V_r}\le \varepsilon_D.
% \]
% Then the mixing magnitudes satisfy
% \[
% \mu_E \le \varepsilon_E,
% \qquad
% \nu_D \le \varepsilon_D,
% \]
% and the perturbation blocks obey
% \begin{align*}
% \op{C_{11}}
% &\le
% \varepsilon_E\,\op{\Sigma_r'}\,\varepsilon_D,
% \\
% \op{C_{12}}
% &\le
% \varepsilon_E\,\op{\Sigma_r'}\,\op{D},
% \\
% \op{C_{21}}
% &\le
% \op{E}\,\op{\Sigma_r'}\,\varepsilon_D,
% \\
% \op{C_{22}}
% &\le
% \op{E}\,\op{\Sigma_r'}\,\op{D}.
% \end{align*}
% In particular:
% \begin{enumerate}[label=(\roman*),leftmargin=2em]
%     \item the direct perturbation of the preserved leading block \(C_{11}\) is second-order in the two mixing levels \(\varepsilon_E,\varepsilon_D\);
%     \item the cross-block couplings \(C_{12}\) and \(C_{21}\) are each first-order in one of the two mixing levels;
%     \item the full perturbation remains norm-bounded by \eqref{eq:global-perturb-bound}.
% \end{enumerate}
% \end{theorem}

% \begin{proof}
% Since \(\op{X}\le \fro{X}\) for any matrix \(X\), we have
% \[
% \mu_E=\op{U_R^\top E U_r}\le \fro{U_R^\top E U_r}\le \varepsilon_E,
% \]
% and likewise \(\nu_D\le \varepsilon_D\). The rest follows immediately from Proposition~\ref{prop:block-bounds}.
% \end{proof}
\begin{theorem}[Soft mixing regularization yields approximate tail confinement]
\label{thm:soft-mix}
Consider the training objective
\[
\mathcal{L}_{\mathrm{total}}
=
\mathcal{L}_{\mathrm{task}}
+
\lambda_E^{\mathrm{mix}}\fro{U_R^\top E U_r}^2
+
\lambda_D^{\mathrm{mix}}\fro{V_R^\top D V_r}^2.
\]
Suppose that at some iterate (or local optimum) one has
\[
\fro{U_R^\top E U_r}\le \varepsilon_E,
\qquad
\fro{V_R^\top D V_r}\le \varepsilon_D.
\]
Then
\[
\mu_E := \|U_R^\top E U_r\|_2 \le \varepsilon_E,
\qquad
\nu_D := \|V_R^\top D V_r\|_2 \le \varepsilon_D.
\]
Consequently, by Proposition~\ref{prop:tail-confinement},
\[
\|C-C_{\mathrm{tail}}\|_2
\le
\|\Sigma_r'\|_2
\sqrt{
(\varepsilon_E\varepsilon_D)^2
+
(\varepsilon_E\|D\|_2)^2
+
(\|E\|_2\varepsilon_D)^2
}.
\]
Hence soft regularization of the mixing terms drives the perturbation toward a pure tail-block perturbation, without requiring the full perturbation norm \(\|C\|_2\) to vanish.
\end{theorem}

\begin{proof}
The inequalities
\[
\mu_E \le \fro{U_R^\top E U_r}\le \varepsilon_E,
\qquad
\nu_D \le \fro{V_R^\top D V_r}\le \varepsilon_D
\]
follow from \(\|X\|_2\le \|X\|_F\). The stated bound then follows immediately from Proposition~\ref{prop:tail-confinement}(iii).
\end{proof}

\begin{remark}[Interpretation]
Theorem~\ref{thm:soft-mix} formalizes the intended preservation--adaptation trade-off:
\begin{itemize}[leftmargin=2em]
    \item the global perturbation size remains controlled by \(\op{E}\op{\Sigma_r'}\op{D}\);
    \item the extent to which this perturbation leaves the pretrained tail block is controlled separately by the mixing terms \(\mu_E\) and \(\nu_D\), or by their regularized surrogates \(\varepsilon_E,\varepsilon_D\).
\end{itemize}
Thus soft regularization does not force the full perturbation to vanish, but it makes leakage into the preserved subspace quantitatively small; equivalently, it makes the off-tail component \(C-C_{\mathrm{tail}}\) quantitatively small, thereby promoting approximate confinement of the adapted update to the pretrained tail singular block.
\end{remark}

\section{Subspace Perturbation Corollary}
Let the singular values of \(W^{\mathrm{pre}}\) satisfy a spectral gap
\[
\gamma := \sigma_R(W^{\mathrm{pre}}) - \sigma_{R+1}(W^{\mathrm{pre}}) > 0.
\]
Then one expects the leading singular subspaces of \(\widetilde W\) to remain close to those of \(W^{\mathrm{pre}}\) whenever the perturbation is small relative to \(\gamma\).

\begin{corollary}[Gap-dependent approximate preservation]
\label{cor:subspace}
Assume \(\gamma>0\) and let \(\widetilde U_R,\widetilde V_R\) denote the leading \(R\)-dimensional left and right singular subspaces of \(\widetilde W = W^{\mathrm{pre}} + \Delta W\). If
\[
\op{\Delta W} < \gamma,
\]
then the principal angle deviations\footnote{The quantity \(\|\sin\Theta(A,B)\|_2\) is a standard subspace-distance measure, equal to the largest sine of the principal angles between the two subspace.} satisfy the standard perturbative estimate
\[
\op{\sinTheta(\widetilde U_R,U_R)}
\;\lesssim\;
\frac{\op{\Delta W}}{\gamma},
\qquad
\op{\sinTheta(\widetilde V_R,V_R)}
\;\lesssim\;
\frac{\op{\Delta W}}{\gamma}.
\]
Using \eqref{eq:global-perturb-bound},
\[
\op{\sinTheta(\widetilde U_R,U_R)},
\ \op{\sinTheta(\widetilde V_R,V_R)}
\;\lesssim\;
\frac{\op{E}\,\op{\Sigma_r'}\,\op{D}}{\gamma}.
\]
Moreover, the leading-block contamination is sharpened by
\[
\op{C_{11}}
\le
\mu_E\,\op{\Sigma_r'}\,\nu_D,
\]
so small \(\mu_E,\nu_D\) imply that the perturbation of the leading block is much smaller than the total perturbation in operator norm.
\end{corollary}

\begin{remark}
Corollary~\ref{cor:subspace} shows the distinction between:
\begin{itemize}[leftmargin=2em]
    \item \emph{global perturbation control}: controlled by \(\op{E}\op{\Sigma_r'}\op{D}\),
    \item \emph{leading-subspace preservation}: improved when \(\mu_E,\nu_D\) are small.
\end{itemize}
Thus \(\mu_E,\nu_D\) are the right quantities to regularize if the theoretical goal is not merely bounded perturbation, but approximate preservation of the pretrained dominant spectral geometry.
\end{remark}

\section{Forward Hidden-State Perturbation}

If a hidden state \(h\in\R^n\) is transformed via the adapted weight \(\widetilde W\), then the update contribution is
\[
\Delta h := \Delta W\,h.
\]
From \eqref{eq:global-perturb-bound},
\[
\norm{\Delta h}_2
\le
\op{\Delta W}\,\norm{h}_2
\le
\op{E}\,\op{\Sigma_r'}\,\op{D}\,\norm{h}_2.
\]
Hence the hidden-state perturbation is always norm-bounded, independently of whether exact block separation holds. Soft mixing regularization does not primarily shrink \(\op{\Delta W}\); rather, it shrinks the \emph{part of \(\Delta W\)} that leaks into or couples with the preserved leading singular structure.

\section{Suggested Ablation Study}

We recommend three training variants:

\begin{enumerate}[leftmargin=2em]
    \item \textbf{No regularizer}: \(\lambda_E^{\mathrm{mix}}=0,\ \lambda_D^{\mathrm{mix}}=0\).
    \item \textbf{\(\mu_E\)-only regularizer}: \(\lambda_E^{\mathrm{mix}}>0,\ \lambda_D^{\mathrm{mix}}=0\).
    \item \textbf{\(\mu_E+\nu_D\) regularizer}: \(\lambda_E^{\mathrm{mix}}>0,\ \lambda_D^{\mathrm{mix}}>0\).
\end{enumerate}

The metrics below are designed to distinguish:
\begin{itemize}[leftmargin=2em]
    \item downstream task performance,
    \item total perturbation size,
    \item leakage into the preserved leading block,
    \item cross-subspace coupling,
    \item actual movement of the leading singular subspaces.
\end{itemize}

\subsection{Recommended reported metrics}

For each adapted layer (or averaged across layers), report:

\begin{enumerate}[label=\arabic*.,leftmargin=2em]
    \item \textbf{Task metric}:
    \[
    \text{TaskScore}
    \in
    \{\text{Accuracy, F1, Exact Match, Rouge-L, BLEU, perplexity, etc.}\}.
    \]
    Choose the standard metric for the downstream benchmark.

    \item \textbf{Total update magnitude}:
    \[
    \mathrm{RelPert}
    :=
    \frac{\op{\Delta W}}{\op{W^{\mathrm{pre}}}}
    \quad\text{or}\quad
    \frac{\fro{\Delta W}}{\fro{W^{\mathrm{pre}}}}.
    \]

    \item \textbf{Left/right mixing magnitudes}:
    \[
    \mu_E = \op{U_R^\top E U_r},
    \qquad
    \nu_D = \op{V_R^\top D V_r}.
    \]

    \item \textbf{Leading-block leakage}:
    \[
    \mathrm{Leak}_{11}
    :=
    \frac{\op{P_{U_R}\Delta W P_{V_R}}}{\op{\Delta W}}
    =
    \frac{\op{C_{11}}}{\op{\Delta W}},
    \]
    where \(P_{U_R}=U_RU_R^\top\) and \(P_{V_R}=V_RV_R^\top\).

    \item \textbf{Cross-block couplings}:
    \[
    \mathrm{Leak}_{12}
    :=
    \frac{\op{P_{U_R}\Delta W (I-P_{V_R})}}{\op{\Delta W}}
    =
    \frac{\op{C_{12}}}{\op{\Delta W}},
    \]
    \[
    \mathrm{Leak}_{21}
    :=
    \frac{\op{(I-P_{U_R})\Delta W P_{V_R}}}{\op{\Delta W}}
    =
    \frac{\op{C_{21}}}{\op{\Delta W}}.
    \]

    \item \textbf{Leading subspace drift}:
    \[
    \mathrm{Drift}_U := \op{\sinTheta(\widetilde U_R,U_R)},
    \qquad
    \mathrm{Drift}_V := \op{\sinTheta(\widetilde V_R,V_R)}.
    \]

    \item \textbf{Leading singular value drift}:
    \[
    \mathrm{SVDrift}
    :=
    \max_{1\le i\le R}
    \frac{\left|\sigma_i(\widetilde W)-\sigma_i(W^{\mathrm{pre}})\right|}
    {\sigma_i(W^{\mathrm{pre}})}.
    \]

    \item \textbf{Efficiency metrics}:
    \[
    \text{TrainableParams},\quad \text{GPU Memory},\quad \text{Train Time / Epoch}.
    \]
\end{enumerate}


\subsection{What to expect empirically}

A useful hypothesis to test is:
\begin{itemize}[leftmargin=2em]
    \item \textbf{No regularizer}: best flexibility, but potentially larger leakage and larger leading-subspace drift.
    \item \textbf{\(\mu_E\)-only}: reduced left-side leakage, but possibly residual right-side coupling through \(D\).
    \item \textbf{\(\mu_E+\nu_D\)}: best preservation of the pretrained leading spectral structure, though possibly at some cost in task performance if the task benefits from repurposing leading pretrained directions.
\end{itemize}

\section{Takeaway}

The key message is:

\begin{center}
\fbox{
\parbox{0.93\linewidth}{
Softly regularizing the singular-basis mixing terms
\(
U_R^\top E U_r
\)
and
\(
V_R^\top D V_r
\)
does \emph{not} merely reduce the total perturbation magnitude; rather, it specifically reduces the part of the perturbation that contaminates or couples with the preserved leading singular structure.
}}
\end{center}

This yields a sharper and more faithful theoretical justification than a global operator-norm bound alone.

\end{document}